\section*{Bab \ref{ch:b1:geological-maps} --- Peta geologi}

\begin{solution}{exo:b1:geological-maps:1}
\omterm{def:b1:geological-maps:attitude}{Jurus}: \omterm{def:b1:spherical-astronomy:horizontal}{azimut} garis horizontal bidang; \omterm{def:b1:geological-maps:attitude}{kemiringan}: sudut antara bidang dan horizontal, diukur tegak lurus \omterm{def:b1:geological-maps:attitude}{jurus}; \omterm{def:b1:geological-maps:attitude}{arah kemiringan}: \omterm{def:b1:spherical-astronomy:horizontal}{azimut} garis paling curam yang menurun. Untuk $135/40$, \omterm{def:b1:geological-maps:attitude}{jurus} menurut kaidah tangan kanan adalah $135 - 90 = \ang{045}$. Simbol: garis \omterm{def:b1:geological-maps:attitude}{jurus} berorientasi
U\,45\,T, dengan tanda pada sisi tenggaranya (ke arah \ang{135}) yang diberi label 40.
\end{solution}

\begin{solution}{exo:b1:geological-maps:2}
$z = 500 - \tan\ang{30}\,(x\sin\ang{90} + y\cos\ang{90}) = 500 - 0.577\,x$.
Pada $(1000, 0)$: $500 - 577 = \qty{-77}{m}$. Pada $(0, 1000)$: \qty{500}{m}, karena bidang itu horizontal sepanjang \omterm{def:b1:geological-maps:attitude}{jurusnya} yang berarah utara--selatan.
\end{solution}

\begin{solution}{exo:b1:geological-maps:3}
Jejak yang mengikuti kontur adalah \omterm{def:b1:geological-maps:geological-map}{kontak} horizontal, pada \qty{350}{m}. Jejak yang berjalan lurus melintasi bukit dan lembah termasuk \omterm{def:b1:geological-maps:geological-map}{kontak} vertikal.
\end{solution}

\begin{solution}{exo:b1:geological-maps:4}
Sepanjang sungai, inti \omterm{def:b1:geological-maps:fold}{antiklin} memperlihatkan sekis dan granit (di bawah batu pasir); inti \omterm{def:b1:geological-maps:fold}{sinklin} ditempati oleh kapur. Satuan tertua berada paling bawah dalam urutan: di tempat lapisan melengkung ke atas dan erosi memotong paling dalam, satuan tertua muncul; di tempat lapisan melengkung ke bawah, satuan termuda terawetkan di tengah.
\end{solution}

\begin{solution}{exo:b1:geological-maps:5}
Dua titik pertama berada pada elevasi yang sama: garis $y = 0$ adalah garis \omterm{def:b1:geological-maps:attitude}{jurus}, sehingga \omterm{def:b1:geological-maps:attitude}{jurusnya} timur--barat. \omterm{def:b1:geological-maps:geological-map}{Kontak} itu \qty{100}{m} lebih rendah pada jarak \qty{500}{m} ke utara: ia miring ke utara, $\tan\delta = 100/500 = 0.2$, $\delta \approx \ang{11.3}$; \omterm{def:b1:geological-maps:attitude}{arah kemiringan} \ang{000}, \omterm{def:b1:geological-maps:attitude}{jurus} \ang{270} menurut kaidah tangan kanan. Dengan teorema: $\overrightarrow{P_1P_2} = (1000, 0, 0)$,
$\overrightarrow{P_1P_3} = (0, 500, -100)$, $\vec n = (0, 100\,000,
500\,000)$, $\cos\delta = 500\,000/509\,902 \approx 0.981$, $\delta \approx
\ang{11.3}$, dan $(n_x, n_y) = (0, 100\,000)$ menunjuk ke utara: hasilnya sama.
\end{solution}

\begin{solution}{exo:b1:geological-maps:6}
Dasar lembah menurun ke utara, $s = 0.03$ (\ang{1.72}). (a) $a = 0 < s$: V menunjuk ke hulu (selatan), mengikuti kontur. (b) $a = -\tan\ang{20} < 0 < s$: V menunjuk ke hulu. (c) $a = \tan\ang{1} \approx 0.017 < s$: V menunjuk ke hulu. (d) $a = \tan\ang{20}
\approx 0.364 > s$: V menunjuk ke hilir (utara).
\end{solution}

\begin{solution}{exo:b1:geological-maps:7}
$\xi_{\text{dasar}} - \xi_{\text{sisi}} = h/(a - s) = 80/(0.364 - 0.03)
\approx \qty{240}{m}$: jejak di sisi lembah berada sekitar \qty{240}{m} ke hulu dari puncak V.
\end{solution}

\begin{solution}{exo:b1:geological-maps:8}
Penutupan \omterm{def:b1:geological-maps:fold}{antiklin} menunjuk ke arah penunjaman: sumbunya menunjam ke utara. Dengan sumbu yang menunjam ke utara, \omterm{def:b1:geological-maps:geological-map}{singkapan} \omterm{def:b1:geological-maps:fold}{sinklin} menutup ke selatan.
\end{solution}

\begin{solution}{exo:b1:geological-maps:9}
Di sisi tenggara-timur tampak satuan yang lebih muda (kapur dan serpih melawan satuan yang lebih tua di seberangnya): blok itu bergerak turun. Tanda-tanda pada simbol sesar berada di sisi itu, sisi yang turun.
\end{solution}

\begin{solution}{exo:b1:geological-maps:10}
Garis yang sama bagi kedua bidang tegak lurus pada kedua normal, sehingga sejajar dengan $\vec n_1 \times \vec n_2$. Untuk $090/30$: $\vec n_1 = (0.577, 0,
1)$; untuk $270/30$: $\vec n_2 = (-0.577, 0, 1)$. Maka $\vec n_1 \times \vec n_2
= (0\cdot 1 - 1\cdot 0,\ 1\cdot(-0.577) - 0.577\cdot 1,\ 0) = (0, -1.155,
0)$: garis utara--selatan horizontal. Bidang pertama miring ke timur, yang kedua ke barat: keduanya menjauhi garis itu, yang merupakan puncak \omterm{def:b1:geological-maps:fold}{antiklin} dengan sumbu horizontal.
\end{solution}

\begin{solution}{exo:b1:geological-maps:11}
Jejak gang vertikal adalah garis lurus yang melalui titik asal dan berjurus \ang{045}, berapa pun elevasinya: di bukit ia berada di suatu tempat pada garis yang sama itu. Gang yang miring adalah bidang $z = -\tan\ang{60}\,d$, dengan $d$ jarak horizontal dari garis \omterm{def:b1:geological-maps:attitude}{jurus} yang melalui titik asal, dihitung positif ke arah tenggara. Pada $z = \qty{300}{m}$: $d = -300/\tan\ang{60}
\approx \qty{-173}{m}$: jejak itu berada \qty{173}{m} di sebelah barat laut garis \omterm{def:b1:geological-maps:attitude}{jurus} yang melalui titik asal, pada sisi ke arah atas \omterm{def:b1:geological-maps:attitude}{kemiringan}.
\end{solution}

\begin{solution}{exo:b1:geological-maps:12}
\[
  \frac{\partial z}{\partial x} \approx \frac{405 - 420}{100} = -0.15, \qquad
  \frac{\partial z}{\partial y} \approx \frac{412 - 420}{100} = -0.08,
\]
sehingga \omterm{def:b1:geological-maps:attitude}{kemiringan} lerengnya adalah $\arctan\sqrt{0.15^2 + 0.08^2} = \arctan 0.17 \approx
\ang{9.6}$. Tanah menurun sepanjang $-\nabla z = (0.15, 0.08)$, yang \omterm{def:b1:spherical-astronomy:horizontal}{azimutnya}
adalah $\arctan(0.15/0.08) \approx \ang{62}$: sel itu menghadap ke timur laut-timur.
\end{solution}

\begin{solution}{pb:b1:geological-maps:1}
\textbf{1.} Sekis, granit, batu pasir, batu gamping, serpih, kapur, aluvium.

\textbf{2.} Dari barat ke timur: kapur dan serpih, batu gamping, batu pasir, lalu
sekis dan granit dekat $x = \qty{3.5}{km}$, batu pasir, batu gamping dan serpih
lagi, kemudian kapur; aluvium di sepanjang seluruh dasar lembah.

\textbf{3.} Batu pasir, dengan sekis dan granit di bawahnya sepanjang sungai: satuan-satuan tertua di inti, sebuah \omterm{def:b1:geological-maps:fold}{antiklin}.

\textbf{4.} Kapur, satuan termuda: sebuah \omterm{def:b1:geological-maps:fold}{sinklin}.

\textbf{5.} Dasar batu pasir paling tinggi pada sumbu \omterm{def:b1:geological-maps:fold}{antiklin}, dan hanya di sepanjang sungai erosi telah menggerus lebih dalam daripadanya.

\textbf{6.} Sebuah ketidakselarasan: batu pasir bertumpu pada dua satuan tua yang berbeda, sekis dan granit, dan memotong garis tepi granit.

\textbf{7.} $435 - 261 = 174$, $432 - 428 = 4$, $512 - 586 = \qty{-74}{m}$.

\textbf{8.} $\overrightarrow{E_1E_2} = (600, 1200, -170)$,
$\overrightarrow{E_1E_3} = (1100, 400, -248)$.

\textbf{9.} Hasil kali silangnya adalah
\[
  (-229\,600,\ -38\,200,\ -1\,080\,000);
\]
diorientasikan ke atas, $\vec n_E = (229\,600,\ 38\,200,\ 1\,080\,000)$.

\textbf{10.} $\|\vec n_E\| \approx 1\,104\,800$; $\cos\delta \approx 0.9776$,
$\delta \approx \ang{12.2}$.

\textbf{11.} \omterm{def:b1:spherical-astronomy:horizontal}{Azimut} $(229\,600, 38\,200)$: $\arctan(229\,600/38\,200)
\approx \ang{80.6}$; \omterm{def:b1:geological-maps:attitude}{orientasi} $081/12$, \omterm{def:b1:geological-maps:attitude}{jurus} $\approx \ang{351}$.

\textbf{12.} Ya: simbol pada sayap timur memiliki garis \omterm{def:b1:geological-maps:attitude}{jurus} kira-kira utara--selatan dan tanda \omterm{def:b1:geological-maps:attitude}{kemiringan} ke timur, berlabel 12.

\textbf{13.} $524 - 669 = -145$, $421 - 459 = -38$, $319 - 209 =
\qty{110}{m}$.

\textbf{14.} $\overrightarrow{W_1W_2} = (700, 1200, 107)$ dan
$\overrightarrow{W_1W_3} = (1100, -600, 255)$. Hasil kali silangnya, yang diorientasikan ke atas, adalah
\[
  (370\,200,\ -60\,800,\ -1\,740\,000) \quad\longrightarrow\quad
  \vec n_W = (-370\,200,\ 60\,800,\ 1\,740\,000).
\]

\textbf{15.} $\cos\delta = 1\,740\,000/\|\vec n_W\| \approx 0.9776$, $\delta
\approx \ang{12.2}$; \omterm{def:b1:geological-maps:attitude}{arah kemiringan}: \omterm{def:b1:spherical-astronomy:horizontal}{azimut} $(-370\,200, 60\,800)$,
$\approx \ang{279.3}$. \omterm{def:b1:geological-maps:attitude}{Orientasi} $279/12$, \omterm{def:b1:geological-maps:attitude}{jurus} $\approx \ang{189}$: sayap barat
miring ke barat-barat laut, sayap timur ke timur-timur laut.

\textbf{16.} $\vec n_E \times \vec n_W = (804 \times 10^6,\ -799\,320 \times
10^6,\ 28\,101 \times 10^6)$ (dibulatkan).

\textbf{17.} Arah menurun adalah $(-804, 799\,320, -28\,101) \times
10^6$: arah $\arctan(-804/799\,320) \approx \ang{0}$, yaitu utara; penunjaman
$\arctan\bigl(28\,101/\sqrt{804^2 + 799\,320^2}\bigr) \approx \ang{2.0}$. Sumbu itu menunjam \ang{2} ke utara.

\textbf{18.} $8000 \times \tan\ang{2.0} \approx \qty{280}{m}$ lebih rendah, hingga sekitar
\qty{-50}{m}. Inti \omterm{def:b1:geological-maps:fold}{antiklin} paling tinggi di dekat sungai dan turun ke utara, sehingga \omterm{def:b1:geological-maps:geological-map}{singkapannya} menutup ke utara, seperti yang diperlihatkan peta.

\textbf{19.} $\overrightarrow{F_1F_2} = (1627, 2500, -276)$ dan
$\overrightarrow{F_1F_3} = (3124, 5500, 77)$ memberikan normal ke atas
\[
  (1\,710\,500,\ -987\,503,\ 1\,138\,500), \qquad
  \delta = \arccos\frac{1\,138\,500}{2\,279\,700} \approx \ang{60},
\]
dan \omterm{def:b1:geological-maps:attitude}{arah kemiringan} $\approx \ang{120}$. \omterm{def:b1:geological-maps:attitude}{Orientasi} $120/60$, \omterm{def:b1:geological-maps:attitude}{jurus}
\ang{030}.

\textbf{20.} Sesar miring ke timur-tenggara, sehingga sisi itu adalah dinding gantung, blok yang berada di atas bidang sesar yang miring; ia bergerak turun (satuan yang lebih muda di permukaan): sesar normal.

\textbf{21.} Bidang sesar miring ke timur-tenggara, ke hilir dan jauh lebih curam daripada dasar sungai menurun: tanah yang lebih rendah bertemu dengan bidang itu lebih jauh ke \omterm{def:b1:geological-maps:attitude}{arah kemiringan}, sehingga jejak membelok ke timur-tenggara di lembah --- aturan-V diterapkan pada bidang sesar.

\textbf{22.} Intrusi granit; erosi sebelum batu pasir; pengendapan kapur; pelipatan; pensesaran; pengendapan aluvium.

\textbf{23.} Setelah 90 juta tahun yang lalu (kapur terlipat) dan sebelum 2.58 juta tahun yang lalu (aluvium tidak tersesarkan); pelipatan terjadi lebih dahulu, pensesaran sesudahnya.

\textbf{24.} Sesar turun \qty{1405}{m} dari F1 (\qty{405}{m}) sampai \qty{-1000}{m}; pada \omterm{def:b1:geological-maps:attitude}{kemiringan} \ang{60} itu memerlukan jarak horizontal sepanjang \omterm{def:b1:geological-maps:attitude}{kemiringan} $1405/\tan\ang{60} \approx \qty{811}{m}$. Sumur harus dibor sekitar \textbf{\qty[text-series-to-math=true]{810}{m} ke timur-tenggara dari F1} (\omterm{def:b1:spherical-astronomy:horizontal}{azimut} \ang{120}).
\end{solution}
